\section{问题三：算力约束下的资源配置优化}
\label{sec:5.4}

\subsection{模型建立}
\label{sec:5.4.1}

给定总预算 $C$（绝对 FLOPs）与外生上下文长度 $L_{ctx}$（token），在参数规模 $N$、训练 token 数 $D$ 与质量 $Q\in[Q_0,1]$ 上最小化问题二的幂次主式。主方案固定问题一在 A4/A5 训练信赖域内选出的 17 域推荐配比 $\bp^*$，这是题目允许的建模选择。记
\begin{equation*}
\begin{aligned}
m_*&=\exp\{s_{\rm red}\tilde\varphi(\bp^*)\}\approx0.8989423,\\
T_N&=AN^{-\alpha}Q^{-\kappa_N},\qquad T_D=BD^{-\beta}Q^{-\kappa_D},\\
L&=E+k_0(1-Q)+m_*(T_N+T_D).
\end{aligned}
\end{equation*}
配比乘子只作用于两项随规模变化的损失，不作用于 $E$ 与质量地板。主质量映射按 P2 取 $Q_{\rm B}=Q_{\rm A}$，故提质起点 $Q_0=0.6624116679870492$。系数 $E,A,\alpha,B,\beta,\kappa_N,\kappa_D,k_0$ 全部读取冻结的 P2 接口，不在本问重拟合。预算约束为
\begin{equation*}
6ND+D\,[g(Q)-g(Q_0)]_++\eta NDL_{ctx}\le C,\qquad Q_0\le Q\le1,\qquad \eta=2\times10^{-4}.
\end{equation*}

作为映射敏感性对照，P2 的锚点映射 $Q_{\rm B}=Q_{\rm A}/Q_{\rm A}(\bp^*)$ 在推荐配比处给出 $Q_0=1$，没有可行提质区间，记为 \texttt{NO\_QUALITY\_RANGE}；若域级质量超出 $[0,1]$ 支持域，按 \texttt{clip\_policy=error} 标记，不作静默截断。

\textbf{提质成本函数。} 需比较题目附录 B 的三类单位 token 提质成本：指数型 $g=10^7e^{6Q}$（凸）、幂函数型 $g=5\times10^9Q^4$（凸）、对数渐近型 $g=2\times10^9\ln(1+10Q)$（凹）。本问按单次训练的静态预算情景建模；题面建议预算 $C=10^{19},10^{22},10^{24}$，硬要求是至少跨越三个数量级，本文自主采用 $10^{19}$ 至 $10^{25}$ 共七档网格并把三个建议档单列。$L_{ctx}$ 是情景输入，不由优化器自由选择；连续 $N,D$ 与 $Q\ge Q_0$ 是为解析求解设定的建模口径。

优化结论继承问题一、二的证据边界：$\bp^*$ 来自 1M 代理模型信赖域，候选收益为预测值；幂次 M4+地板是问题二指定主式，质量项由半合成 B6 的 360 点估计、B7 新增 90 点作插值留出，主式留出 RMSE 约 $0.053$（线性对照约 $0.042$），$\kappa_D$ 较弱。因此最优配置是给定该标度律、映射与参数后的模型内方案，配比乘子通道亦为探索性假设。

\textbf{降维与闭式对照。} 三项成本都与 $D$ 成正比，而 $\partial L/\partial D=-m_*\beta T_D/D<0$，故预算用尽。记 $K=6+\eta L_{ctx}$、$\Delta g=[g(Q)-g(Q_0)]_+$，由约束取等得
\begin{equation*}
D=\frac{C}{KN+\Delta g},
\end{equation*}
问题降为 $(\log_{10}N,Q)$ 的二维有界优化。固定 $Q$、令 $\Delta g=0$ 的解析对照为
\begin{equation*}
N^*=\left[\frac{\alpha A Q^{-\kappa_N}}{\beta B Q^{-\kappa_D}}\right]^{1/(\alpha+\beta)}
\left(\frac{C}{K}\right)^{\beta/(\alpha+\beta)},\qquad D^*=\frac{C}{KN^*},
\end{equation*}
式中 $m_*$ 与地板因对固定 $Q$ 求规模导数而抵消，但仍影响不同 $Q$ 之间的选择。由此
\begin{equation*}
\frac{D^*}{N^*}\propto\left(\frac{C}{K}\right)^{(\alpha-\beta)/(\alpha+\beta)}
Q^{2(\kappa_N-\kappa_D)/(\alpha+\beta)},
\end{equation*}
两个指数在当前参数下约为 $0.09677$ 与 $0.42065$；这是无提质开销对照，真实 $\Delta g>0$ 情景由数值解给出。

上下文长度的解析临界值由注意力项与训练项之比 $\eta L_{ctx}/6$ 给出：
\begin{equation*}
L_{ctx}^{\rm crit}=\frac{6}{\eta}=30000\ \text{token},
\end{equation*}
在固定 $L_{ctx}$ 下该比值与 $C,N,D$ 无关；C7 五档中 8192 在临界值以下、32768 在其以上，30000 本身不是 C7 实测档位。

\textbf{结构性转移的定义与判据。} 本文把最优策略的质变操作化为全局最优解的质量约束活跃状态变化：\texttt{LOWER}（$Q^*=Q_0$）、\texttt{INTERIOR}、\texttt{UPPER}（$Q^*=1$）；若两个内点局部极小分支交换全局最优地位，即使同属 \texttt{INTERIOR} 也单列全局分支切换。这一做法有直接的文献先例：Prescriptive Scaling Laws\cite{lovelace} 在数据受限前沿上发现 turn-back 现象——预算越过临界后最优策略由「加重复次数」质变为「加大模型、减轮次」。对固定 $Q$ 的规模最优解，令 $S=KN+\Delta g$，剖面目标的质量导数与边际比为
\begin{equation*}
F_Q=-k_0-\frac{m_*}{Q}(\kappa_NT_N+\kappa_DT_D)+\frac{m_*\beta T_D\,g'(Q)}{S},\qquad
\Phi(C;Q)=\frac{\left[k_0+\dfrac{m_*}{Q}(\kappa_NT_N+\kappa_DT_D)\right]S}{m_*\beta T_D\,g'(Q)}.
\end{equation*}
其中 $N,D$ 均在该固定 $Q$ 下优化。内层光滑且未触边界时，$\Phi(C;Q_0)\le1$、内点 $\Phi(C;Q)=1$、$\Phi(C;1)\ge1$ 分别为相应状态的局部一阶必要条件；$\Phi=1$ 的根不是全局切换的充分条件，须分别输出局部根与全局临界点。$Q_0\ge1-10^{-10}$ 时记 \texttt{NO\_QUALITY\_RANGE}，搜索区间未括根记 \texttt{NOT\_BRACKETED}。公式正确性已作单点有限差分核验：$C=10^{20}$、$L_{ctx}=2048$、指数型、$Q=0.8$ 处解析导数 $-0.1164515205405$ 与差分 $-0.1164515204932$ 的相对误差约 $4.1\times10^{-10}$（对应 $\Phi=1.63892372$）。

\textbf{敏感性与验收。} 敏感性覆盖三类 $g$、C7 五档 $L_{ctx}$、$\eta=10^{-4},2\times10^{-4},3\times10^{-4}$、P2 参数重抽样、主/锚点两种质量映射、P1 域质量与配比不确定性，以及固定配比乘子 \texttt{m=1}、H$_{\rm tot}$ 与无地板对照。固定 $\bp^*$ 时 H$_{\rm tot}$ 对完整基线乘一个与 $N,D,Q$ 无关的正常数，故其最优 $N,D,Q$ 与 \texttt{m=1} 相同、最优 Loss 按该常数缩放。验收阈值：预算超限与三项份额和误差均不大于 $10^{-9}$，预算使用率与 1 之差小于 $10^{-9}$，无提质时闭式 $N^*$ 误差小于 $10^{-5}$，关键 18 组求解 Loss 与独立粗网格之差不超过 $2\times10^{-5}$。\texttt{P3\_optimal\_config.csv} 的份额分母统一为实际总消耗，同时另报预算使用率。

\subsection{求解方法}
\label{sec:5.4.2}

消去 $D$ 后，在 $(\log_{10}N,Q)$ 上求解：固定 $Q$ 的内层对 $\log_{10}N$ 作有界一维搜索并核对边界未截住解，外层对 $Q$ 加密搜索、保护 $Q_0$ 与 $1$ 两个端点，同时比较多个内点候选与独立粗网格。局部 $\Phi=1$ 根与全局状态切换分别识别，仅当区间已括根且确认为单次切换时才用二分定位临界预算。主网格为 3 种成本函数 $\times$ 7 档预算 $\times$ 5 档上下文长度 $\times$ 2 个起始质量对照，共 210 组；关键 18 组同时用粗网格与多起点 SLSQP 核对。对 $Q_0=1$、未括根情形分别记录 \texttt{NO\_QUALITY\_RANGE}、\texttt{NOT\_BRACKETED} 与空数值。本问已按幂次主式重跑并生成运行 manifest（\texttt{P3\_run\_manifest.json}，北京时间 2026-09-26 11:37，含 P1/P2/C7 输入 SHA-256）；\texttt{q3\_03\_mixture\_joint.py} 改为运行即退出，不作为主结论。

\subsection{结果与分析}
\label{sec:5.4.3}

\textbf{结果证据范围。} 目标式已固定为问题二的幂次 M4+质量地板并保留配比乘子 $m_*\approx0.8989423$，$Q_0=0.6624116679870492$。\texttt{src/code\_q3/} 已按该口径重跑（manifest 记录北京时间 2026-09-26）。主网格为 3 种成本函数 $\times$ 7 档预算 $\times$ 5 档 C7 上下文长度 $\times$ 2 个同映射下的起始质量对照，共 210 组；题面建议的三档预算 $10^{19},10^{22},10^{24}$ FLOPs 在结果中单列。标度律质量项拟合于 B6（\texttt{supplementary\_NQ\_experiment.csv}，360 点，$6ND$ 预算 $\log_{10}C\in[18.62,22.63]$、$N\le1.2\times10^{10}$、$D\le6\times10^{11}$），B7 同域留出，B8 将观测扩展至 $\log_{10}C\approx24.92$（$N\le7\times10^{11}$、$D\le2\times10^{12}$）。据此对每个最优解按其 $N^*,D^*$ 相对支撑范围标注三级可靠性 \texttt{IN\_FIT}／\texttt{IN\_OBS}／\texttt{EXTRAP}（见\textbf{高预算外推警告}）；关键结论是\textbf{所有临界预算均落在 B6/B7 拟合域内}，而 $\log_{10}C\gtrsim23.5$ 的高预算最优规模配置为纯外推。质量参数来自半合成 B6、配比通道为探索性，证据边界与问题二一致。下文先给出可由主式推出的解析关系，再给出数值最优配置、成本份额、配置比斜率诊断与临界预算。

\textbf{上下文长度临界值。} 注意力项与训练项之比恒为 $\eta L_{ctx}/6$，二者相等给出 $L_{ctx}^{\rm crit}=30000$ token，与预算、成本函数和规模无关；8192 在临界值以下、32768 在其以上。注意力份额随 $L_{ctx}$ 增大而上升，$L_{ctx}=32768$ 时已接近或超过训练份额，越过临界值后最优模型随之缩小。成本交叉点与预算驱动的策略转移是两件事，须分开陈述。

\textbf{数值最优配置。} 表~\ref{tab:q3opt} 给出 $L_{ctx}=2048$、主口径下三档建议预算与三类提质成本对应的最优配置。低预算 $C=10^{19}$ 时，指数型与幂函数型成本下最优质量停留在下界（\texttt{LOWER}，不提质），对数型成本因提质边际成本低而直接提满（\texttt{UPPER}）；到 $C=10^{22}$、$10^{24}$，三种成本函数下最优解均为 $Q^*=1$，最优 $N,D$ 沿算力前沿增长、最优 Loss 稳步下降。$10^{22}$ 及以上超出 B6 训练覆盖，属外推结果。

\begin{table}[H]
\centering\small
\caption{三档建议预算下的最优资源配置（$L_{ctx}=2048$，主口径；数据源 \texttt{P3\_optimal\_config.csv}，2026-09-26 11:37）}
\label{tab:q3opt}
\begin{tabularx}{\textwidth}{@{}l c c c c c c@{}}
\toprule
成本 $g(Q)$ & $\log_{10}C$ & $Q^*$ & $N$ & $D$ & 最优 Loss & 状态 \\
\midrule
指数型 & 19 & 0.6624 & $2.414\times10^{8}$ & $6.464\times10^{9}$ & 2.9780 & LOWER \\
指数型 & 22 & 1 & $5.584\times10^{9}$ & $2.545\times10^{11}$ & 2.1038 & UPPER \\
指数型 & 24 & 1 & $4.072\times10^{10}$ & $3.780\times10^{12}$ & 1.8914 & UPPER \\
幂函数型 & 19 & 0.6624 & $2.414\times10^{8}$ & $6.464\times10^{9}$ & 2.9780 & LOWER \\
幂函数型 & 22 & 1 & $5.663\times10^{9}$ & $2.479\times10^{11}$ & 2.1046 & UPPER \\
幂函数型 & 24 & 1 & $4.082\times10^{10}$ & $3.764\times10^{12}$ & 1.8914 & UPPER \\
对数型 & 19 & 1 & $3.167\times10^{8}$ & $3.620\times10^{9}$ & 2.9380 & UPPER \\
对数型 & 22 & 1 & $5.140\times10^{9}$ & $2.969\times10^{11}$ & 2.0988 & UPPER \\
对数型 & 24 & 1 & $4.023\times10^{10}$ & $3.867\times10^{12}$ & 1.8910 & UPPER \\
\bottomrule
\end{tabularx}
\end{table}

\textbf{成本份额。} 表~\ref{tab:q3share} 给出三档预算的训练、注意力与提质份额（分母为实际总消耗，预算使用率核验为 1）。$C=10^{19}$ 下训练成本占主导（指数型约 $93.6\%$）；对数型成本在 $C=10^{19}$ 已把约 $26.5\%$ 投入提质，与其低边际成本一致。注意力份额随 $L_{ctx}$ 显著上升，是 $L_{ctx}^{\rm crit}=30000$ 的直接体现。

\begin{table}[H]
\centering\small
\caption{三档预算下的成本份额（$L_{ctx}=2048$，分母为实际总消耗；数据源 \texttt{P3\_optimal\_config.csv}）}
\label{tab:q3share}
\begin{tabularx}{\textwidth}{@{}l c c c c@{}}
\toprule
成本 $g(Q)$ & $\log_{10}C$ & 训练份额 & 注意力份额 & 提质份额 \\
\midrule
指数型 & 19 & 93.61\% & 6.39\% & 0 \\
指数型 & 22 & 85.27\% & 5.82\% & 8.91\% \\
指数型 & 24 & 92.37\% & 6.31\% & 1.32\% \\
幂函数型 & 19 & 93.61\% & 6.39\% & 0 \\
幂函数型 & 22 & 84.24\% & 5.75\% & 10.01\% \\
幂函数型 & 24 & 92.19\% & 6.29\% & 1.52\% \\
对数型 & 19 & 68.77\% & 4.69\% & 26.54\% \\
对数型 & 22 & 91.57\% & 6.25\% & 2.18\% \\
对数型 & 24 & 93.34\% & 6.37\% & 0.28\% \\
\bottomrule
\end{tabularx}
\end{table}

\textbf{结构性转移与临界预算。} $L_{ctx}=2048$ 主口径下，质量状态随预算变化：指数型在 $\log_{10}C\approx19.059$ 离开下界、$\approx20.355$ 到达 $Q=1$；幂函数型为 $19.529$ 与 $20.106$；对数型离开下界与到达上界同发生于 $\approx18.579$（局部 $\Phi=1$ 根约 $18.904$，与全局切换相差 $0.325$ dex）。P2 条件 bootstrap（300 次全部括根）给出离开下界的 $5\%$–$95\%$ 区间：指数型 $[19.038,19.084]$、幂函数型 $[19.503,19.560]$、对数型 $[18.539,18.620]$。由此得到清晰的预算阈值规律：对数型成本阈值最低、指数型居中、幂函数型最高，提质是否划算随成本函数形状有序移动。

\textbf{P1 不确定性传播（Q3-9）。} 进一步将问题一的域质量 $Q_i$ 与配比 $\bp^*$ 的 bootstrap 样本（各 300 次）逐样本重算 $Q_0(\bp^*)=\sum_i p_i^* Q_i$ 并与 P2 参数样本联合，向量化重新识别临界预算，300 次全部括根。加权起始质量本身的抽样波动很小，$Q_0(\bp^*)=0.66237\pm0.00203$（5\%–95\% 区间 $[0.65904,0.66553]$）。P1+P2 联合的离开下界 5\%–95\% 区间为指数型 $[19.031,19.091]$、幂函数型 $[19.501,19.563]$、对数型 $[18.540,18.621]$，到达 $Q=1$ 为指数型 $[20.292,20.428]$、幂函数型 $[20.039,20.179]$、对数型 $[18.540,18.621]$。与仅含 P2 的区间相比，加入 P1 通道仅使区间边界外扩约 $0.003$ dex（例如指数型离开下界由 $[19.038,19.084]$ 微展至 $[19.031,19.091]$），三种成本函数的阈值排序与"提质是否划算"的结论对配比／域质量抽样稳健。

\textbf{求解核验。} 本轮预算残差 $3.93\times10^{-16}$、三项份额和误差 $3.33\times10^{-16}$，均小于 $10^{-9}$；无提质开销时闭式 $N^*$ 相对误差 $1.54\times10^{-7}$（阈值 $10^{-5}$）；关键组求解 Loss 与独立粗网格最大差 $1.80\times10^{-5}$（阈值 $2\times10^{-5}$）；搜索边界未截住主解。锚点映射 $Q_B=Q_A/Q_A(\bp^*)$ 下 $Q_0=1$，记为 \texttt{NO\_QUALITY\_RANGE}；\texttt{pile\_cc} 单域映射质量越界记为 \texttt{OUT\_OF\_SUPPORT} 且不截断。

\textbf{配置比斜率诊断（Q3-10）。} 由主式可得无提质开销时最优规模比满足幂律 $D^*/N^*\propto (C/K)^{(\alpha-\beta)/(\alpha+\beta)}$，即 $\mathrm{d}\ln(D^*/N^*)/\mathrm{d}\ln C$ 的理论基线为 $(\alpha-\beta)/(\alpha+\beta)=0.09677$（$\alpha=0.34,\beta=0.28$）。逐预算数值微分证实：在低预算质量锁定于 $Q_0$ 的区段，斜率与基线偏离仅 $\sim10^{-9}$；而当预算跨过提质阈值、$Q^*$ 快速上行挤占数据预算时，该斜率出现显著负向拐点。主口径 $L_{ctx}=2048$ 下偏离绝对值最大处（拐点）为对数型 $\log_{10}C\approx18.55$、幂函数型 $\approx19.60$、指数型 $\approx19.10$，与前述质量约束切换预算（对数型 $18.58$、幂函数型 $19.53$、指数型 $19.06$）一致，二者互为印证：结构性转移既体现在质量约束状态的切换，也体现在规模配置比增速的可观测拐折。诊断明细见 \texttt{P3\_dn\_slope\_diagnostic.csv} 与 \texttt{P3\_dn\_slope\_kink.csv}。

\textbf{高预算外推警告。} 标度律质量项拟合于 B6（$\log_{10}C\in[18.62,22.63]$、$N\le1.2\times10^{10}$、$D\le6\times10^{11}$），B7 同域留出，B8 将观测扩展至 $\log_{10}C\approx24.92$。按每个最优解的 $N^*,D^*$ 相对该支撑标注：题面三档中 $10^{19},10^{22}$ FLOPs 及全部临界预算（离开下界 $18.58$–$19.53$、到达 $Q=1$ $18.58$–$20.36$）均落在 B6/B7 拟合域内，故\textbf{结构性转移与提质阈值结论非外推}；纯外推（\texttt{EXTRAP}）起点约 $\log_{10}C\approx23.5$（主口径 $L_{ctx}=2048$，指数/幂函数型 $23.55$、对数型 $23.5$），$10^{24}$ 档已进入外推区。每条成本函数在主网格 7 档预算 $\times5$ 档上下文共 35 组中，\texttt{IN\_FIT}／\texttt{IN\_OBS}／\texttt{EXTRAP} 分别为 $21/5/9$；因此 $10^{24}$ 及以上的最优规模配置须谨慎解读。标注明细见 \texttt{P3\_extrapolation\_by\_config.csv} 与 \texttt{P3\_extrapolation\_warning.json}。

\textbf{证据闭合说明（问题三）。} 问题一已提供域质量 $Q_i$ 与配比 $\bp^*$ 的逐次 bootstrap 接口（\texttt{P1\_bootstrap\_Q\_domain.npz}、\texttt{P1\_bootstrap\_p\_star.npz}），Q3-9 已据此完成固定配比（仅 P2）与重选配比（P1+P2）两类不确定性的分别传播，全链条区间已闭合（见上"P1 不确定性传播"）；结论对配比／域质量抽样稳健。质量参数仍来自半合成 B6、配比通道为探索性，此为与问题二一致的数据边界，非方法缺口。
